\documentclass[../../../include/open-logic-section]{subfiles}

\begin{document}

\olfileid{sth}{choice}{hartogs}
\olsection{可比性与Hartogs（哈托格斯）引理}

好的一面就说这些。坏的一面则是，没有选择公理，情况会变得\emph{混乱}。为了看清这一点，这里给出一个由\cite{Hartogs1915}得到的优美结果：

\begin{lem}[\emph{在 $\ZF$ 中}]\ollabel{HartogsLemma}
对于任意集合 $A$，存在序数 $\alpha$ 使得 $\cardnless{\alpha}{A}$。
\end{lem}

\begin{proof}
如果 $B \subseteq A$ 且 $R \subseteq B^2$，那么根据\olref[ord-arithmetic][using-addition]{rankcomputation}，有$\tuple{B, R} \subseteq
V_{\setrank{A}+4}$。因此，运用分离公理，考虑集合：
\[
	C = \Setabs{\tuple{B, R} \in V_{\setrank{A}+5}}{B\subseteq A
	\text{ 且 $\tuple{B, R}$ 是良序}}
\]
运用替换公理和\olref[ordinals][ordtype]{thmOrdinalRepresentation}，构成集合：
\[
	\alpha = \Setabs{\ordtype{B, R}}{\tuple{B, R} \in C}.
\]
根据\olref[ordinals][basic]{corordtransitiveord}，$\alpha$ 是一个序数，因为它是一个由序数构成的传递集。事实上，若 $\gamma \in \beta \in \alpha$，则存在 $B \subseteq R$ 使得 $\beta = \ordtype{B, R}$，于是由\olref[ordinals][iso]{wellordinitialsegment}，存在 $b \in B$ 使得$\gamma = \ordtype{B_b, R_b}$，从而 $\gamma \in \alpha$。

用反证法，假设存在单射 $f \colon \alpha \to A$。令：
\begin{align*}
	B &= \ran{f}\\
	R &= \Setabs{\tuple{f(\alpha), f(\beta)} \in A \times A}{\alpha \in \beta}.
\end{align*}
显然 $\alpha = \ordtype{B, R}$ 且 $\tuple{B, R} \in C$。所以 $\alpha \in \alpha$，矛盾。
\end{proof}

这蕴涵了一个深刻的结果：

\begin{thm}[\emph{在 $\ZF$ 中}]
下列命题等价：
\begin{enumerate}
	\item\ollabel{equivwo} 良序公理
	\item\ollabel{equivcompare} 对任意集合 $A$ 和 $B$，要么 $\cardle{A}{B}$ 要么 $\cardle{B}{A}$
\end{enumerate}
\end{thm}

\begin{proof}
\emph{\olref{equivwo} $\Rightarrow$ \olref{equivcompare}。}固定 $A$ 和 $B$。由\olref{equivwo}，存在良序 $\tuple{A, R}$ 和 $\tuple{B, S}$。根据\olref[ordinals][ordtype]{thmOrdinalRepresentation}，令 $f \colon \alpha \to \tuple{A, R}$ 和 $g \colon \beta \to \tuple{B, S}$ 为同构映射。由\olref[sth][ordinals][basic]{ordinalsaresubsets}，要么 $\alpha \subseteq \beta$ 要么 $\beta \subseteq \alpha$。若 $\alpha \subseteq \beta$，则$\comp{f^{-1}}{g} \colon A \to B$是单射，因而 $\cardle{A}{B}$；类似地，若 $\beta \subseteq \alpha$ 则 $\cardle{B}{A}$。

\emph{\olref{equivcompare} $\Rightarrow$ \olref{equivwo}。}固定 $A$；由\olref{HartogsLemma}，存在某个序数 $\beta$ 使得 $\cardnless{\beta}{A}$。根据\olref{equivcompare}，我们有 $\cardle{A}{\beta}$。所以存在某个单射 $f \colon A \to \beta$，利用此单射，我们可以通过定义序关系$\Setabs{\tuple{a, b} \in A \times A}{f(a)
\in f(b)}$来对 $A$ 的元素进行良序化。
\end{proof}
\noindent
一个直接的推论是：如果良序公理不成立，那么有些集合在大小上是\emph{真正不可比的}。因此，如果良序公理不成立，超限基数运算将会混乱。例如，我们将不得不放弃这样的想法：若 $A$ 和 $B$ 都是无限的，则$\cardeq{\cardeq{A \disjointsum B}{A \times
B}}{M}$，其中 $M$ 是 $A$ 和 $B$ 中较大者（参见\olref[card-arithmetic][simp]{cardplustimesmax}）。问题很简单：如果我们无法\emph{比较} $A$ 和 $B$ 的大小，那么讨论哪个更大便毫无意义。

\end{document}